\section{Introduzione: problema, requisiti e sintesi dei risultati}\label{sec:intro}

\subsection{Il problema}\label{sec:intro-problema}

La traccia di riferimento osserva che «riconoscere un oggetto o un soggetto in movimento non basta: un ostacolo può nasconderlo temporaneamente, mentre la scena continua a cambiare». Chiede un prototipo capace di «rilevare elementi mobili e ostacoli, gestire interruzioni della visibilità e distinguere le informazioni osservate da quelle incerte», che dia «feedback spaziali chiari e tempestivi quando le condizioni lo consentono, segnalando invece la perdita di affidabilità quando necessario», con «continuità percettiva, bassa latenza ed efficienza elaborativa, senza rinunciare alla sicurezza delle indicazioni».

Il riconoscimento non basta perché un rilevatore non produce nulla nell'istante in cui il soggetto esce dalla vista. Da quell'istante un sistema può tacere, e il silenzio è letto come assenza; può continuare a mostrare l'ultima posizione, o una posizione estrapolata, come se fosse osservata; oppure può dichiarare che cosa sa e con quale livello. Le prime due scelte sono indicazioni non sicure: un filtro privo di misure accumula errore per tutta l'occlusione \cite{cao2023}, e in simulazione, con modello ideale del sensore, un punto lasciato fermo all'ultima misura dista dal soggetto $1{,}48$--$1{,}54\unita{m}$ in mediana dopo $2\unita{s}$, contro $0{,}37$--$0{,}39\unita{m}$ del punto previsto (una disposizione, miscela di comportamenti, $N$ fra 2077 e 2115 passaggi) \prova{S}. Questo allegato descrive un sistema che adotta la terza scelta.

Il sistema usa una camera di profondità a stereo attivo, frontale e ferma, ad altezza $h=1{,}2\unita{m}$ e inclinazione $\theta=15^\circ$ verso il basso, con portata utile di $3\unita{m}$. La catena di elaborazione è: mappa polare del suolo a quattro stati (libero, occupato osservato, in ombra, non osservato); gruppi e regola «mobile»; associazione e filtro di Kalman \cite{kalman1960} a velocità quasi costante; durante un'occlusione, una regione di previsione $\Reg_\alpha(t)$ con livello dichiarato e un avviso di «possibile presenza» sull'ombra dell'ostacolo; regole scritte, non apprese; uscita e segnale di «affidabilità persa». Il sistema osserva e avvisa: nessuna uscita è un comando, e non vi sono funzioni d'ingaggio o di puntamento.

\subsection{Requisiti misurabili}\label{sec:intro-requisiti}

La Tabella~\ref{tab:intro-requisiti} traduce ogni requisito della traccia in una grandezza, con il modo di misurarla e la soglia di riferimento. L'unità statistica è il \emph{passaggio}: un soggetto attraversa il campo, scompare dietro un ostacolo fermo e di norma ricompare. Le definizioni complete e gli intervalli stanno in \S\ref{sec:val}; le soglie sono criteri di progetto e il loro esito è nella Tabella~\ref{tab:intro-sintesi}. Ai requisiti della traccia il progetto aggiunge un componente di classificazione su più classi di oggetti in movimento (persone; veicoli di tipo diverso, per esempio auto, treni, aerei a terra; altri ostacoli), ciascuna con un profilo cinematico, e un profilo deterministico, da configurazione iniziale, di chi porta il sensore.

\begin{table}[!t]
\centering\footnotesize
\caption{Requisiti della traccia e grandezze che li misurano. $\tocc$: tempo trascorso dall'ultima osservazione; $\ell^-$: estremo inferiore dell'intervallo di Wilson al 95\,\% \cite{wilson1927}; $\Delta=1/15\unita{s}$: periodo di campionamento.}\label{tab:intro-requisiti}
\begin{tabularx}{\textwidth}{@{}>{\RaggedRight\arraybackslash}p{2.5cm}LL>{\RaggedRight\arraybackslash}p{3.7cm}@{}}
\toprule
Requisito & Grandezza & Come si misura & Soglia di riferimento\\
\midrule
Rilevare elementi mobili & quota dei soggetti a $0{,}5\unita{m/s}$ dichiarati «in movimento» entro $1\unita{s}$; quota dei passaggi in cui un oggetto fermo è dichiarato «in movimento» & conteggio per passaggio, con intervallo & almeno 9 su 10 fino a $2\unita{m}$; la seconda quota è riportata con il suo intervallo\\
Rilevare ostacoli & ostacolo fermo localizzato entro $15\unita{cm}$ dalla posizione vera & conteggio per prova & 5 prove su 5 fra $1{,}1$ e $3\unita{m}$\\
Gestire le interruzioni della visibilità: continuità & quota dei cicli senza osservazione in cui è emessa una regione $\Reg_\alpha(t)\neq\emptyset$; passaggi con almeno un ciclo senza regione & sui passaggi con almeno 4 misure prima della scomparsa & 0 passaggi con un'interruzione\\
\quad copertura & quota dei passaggi ancora nascosti a $\tocc=\tau$ con il soggetto in $\Reg_\alpha$, al livello dichiarato $1-\alpha=0{,}90$; area di $\Reg_\alpha$ rispetto all'ombra & per comportamento e per $\tau\in\{0{,}5;1;2;3\}\unita{s}$ & $\ell^->0{,}90$ fino a $2\unita{s}$ e $\ell^->0{,}85$ a $3\unita{s}$, in ogni combinazione\\
\quad identità & quota dei passaggi con la stessa identità alla ricomparsa & sui passaggi con ricomparsa & almeno $0{,}95$ per chi resta sul proprio percorso\\
Distinguere osservato e incerto & uscite «osservato» emesse mentre il soggetto è nascosto & conteggio sui cicli con il soggetto nascosto & 0\\
Feedback chiari e tempestivi & tempo dal primo fotogramma alla chiamata dell'uscita sonora; presenza di stato, età e livello in ogni uscita & misura di tempo; controllo sullo schema dell'uscita & al più $300\unita{ms}$\\
Segnalare la perdita di affidabilità & ritardo del segnale dedicato; segnali falsi & prova con flusso interrotto, sensore coperto o spostato & entro $1\unita{s}$; 0 falsi in $60\unita{s}$\\
Bassa latenza ed efficienza & latenza $\tau$ dall'evento all'uscita; tempo d'orologio per ciclo su un filo di CPU & 95º percentile; numero di cicli oltre $\Delta$ & 95º percentile di $\tau$ entro $150\unita{ms}$; ciclo entro $\Delta=66{,}7\unita{ms}$\\
Sicurezza delle indicazioni & celle «libero» contenute nell'ombra vera; istanti con il soggetto nascosto fuori dall'area d'avviso; avvisi tolti senza osservazione diretta & conteggi su celle, istanti e passaggi & 0 per ciascuna\\
Classificazione e profili (requisito di progetto) & differenze fra due esecuzioni sugli stessi fotogrammi che differiscono solo nelle etichette di classe & confronto di celle «libero», area d'avviso e livelli & 0 celle «libero» in più; 0 avvisi ridotti; nessun livello più basso\\
\bottomrule
\end{tabularx}
\end{table}

\subsection{L'idea}\label{sec:intro-idea}

Il sistema tiene distinte tre informazioni e non le mescola in nessuna uscita (Figura~\ref{fig:intro-catena}). La prima è ciò che è \emph{osservato} nel fotogramma corrente: il suolo visto libero $\Oss(t)$, gli ostacoli, i soggetti con una misura nel ciclo. La seconda è ciò che \emph{non è osservato ma è previsto}: per un soggetto seguito ed entrato nell'ombra $\Omb(t)$ di un ostacolo, una regione $\Reg_\alpha(t)$ emessa a ogni ciclo con il tempo trascorso $\tocc$ e, quando è dichiarabile, il livello; quando il livello non è dichiarabile la regione resta, con lo stato «previsto ampio». La terza è ciò che \emph{non è noto}: fuori campo, oltre la portata, pixel senza dato, ombra in cui non è entrato alcun soggetto seguito.

Il principio di sicurezza è uno: l'assenza di un rilevamento non è mai un'indicazione di via libera. Lo stato «libero» è scritto solo per il suolo osservato nel fotogramma corrente, e nessun contatore di tempo, modello appreso, etichetta di classe, identità o previsione lo produce. L'insieme di possibile presenza $\Pres(t)$ si riduce solo dove il suolo è rivisto libero; previsione e classe possono allargare un avviso, mai restringerlo. È lo stesso principio dei sistemi industriali di sorveglianza tridimensionale, che trattano come occupato ciò che nessun sensore vede \cite{therobotreport-veo}, e dei metodi che sovrastimano lo stato degli ostacoli nascosti \cite{orzechowski2018}. Quando il sensore non è in salute il sistema emette un segnale dedicato, dichiara tutta la mappa «non osservato» e continua a uscire a ogni ciclo.

\begin{figure}[H]
\centering
\begin{tikzpicture}[font=\scriptsize,>=Stealth,node distance=0.42cm,
  b/.style={draw,rounded corners=1pt,align=center,text width=2.05cm,minimum height=0.95cm,inner sep=2pt}]
\node[b] (cam) {camera di profondità ferma};
\node[b,right=of cam] (map) {mappa polare del suolo a quattro stati};
\node[b,right=of map] (gru) {gruppi e regola «mobile»};
\node[b,right=of gru] (fil) {associazione e filtro di Kalman};
\node[b,right=of fil] (occ) {occlusione: $\Reg_\alpha(t)$ e $\Pres(t)$};
\node[b,right=of occ] (reg) {regole e uscita};
\foreach \a/\c in {cam/map,map/gru,gru/fil,fil/occ,occ/reg} \draw[->] (\a)--(\c);
\node[b,below=0.5cm of cam] (sal) {salute del sensore};
\node[b,dashed,below=0.5cm of fil] (cla) {classe: richiesta, qui simulata};
\node[b,text width=4.55cm,below=0.5cm of $(occ)!0.5!(reg)$] (usc) {mai mescolati: osservato; previsto, con regione, livello e $\tocc$; non noto};
\draw[->] (cam)--(sal);
\draw[->,dashed] (cla)--(fil);
\draw[->] (reg.south)--(reg.south|-usc.north);
\draw[->] (sal.south)--++(0,-0.3)-|node[pos=0.25,below]{«affidabilità persa»: scavalca le regole}(usc.south);
\end{tikzpicture}
\caption{La catena di elaborazione e le tre informazioni in uscita. Il tratteggio indica il componente di classificazione, richiesto dal sistema e in questo lavoro simulato con etichette assegnate a mano.}\label{fig:intro-catena}
\end{figure}

\subsection{Contributi}\label{sec:intro-contributi}

\begin{enumerate}[label=\arabic*.]
\item La geometria d'osservazione in forma chiusa: dove il suolo è visibile, l'ombra di un ostacolo, il rumore di profondità e l'effetto di un errore di posa (\S\ref{sec:geo}).
\item Una mappa polare del suolo a quattro stati in cui «libero» è scritto solo dal fotogramma corrente, con la proposizione di sicurezza, i controesempi alle varianti con memoria e il confronto con le mappe persistenti e a voxel (\S\ref{sec:map}).
\item La catena di rilevamento, associazione e stima, con la prova che la massa non è identificabile da misure di posizione e che l'accelerazione non è stimabile con precisione utile; il componente di classificazione su più classi, con i profili cinematici di classe e i requisiti che il sistema gli pone (\S\ref{sec:sti}).
\item Il trattamento dell'occlusione: insieme di possibile presenza con contenimento e monotonia, fascia di guardia, regione di previsione a livello dichiarato che non si svuota mai, e l'enunciato che previsione ed etichetta non riducono l'avviso (\S\ref{sec:occ}).
\item Un metodo di valutazione definito sul passaggio: misure, intervalli, numerosità, limiti della taratura per quantili, lettura di un conteggio «dentro $k$ su $n$» (\S\ref{sec:val}).
\item Il costo di calcolo e le alternative escluse, ciascuna con il controesempio che la esclude (\S\ref{sec:cos}).
\item Un'uscita che porta sempre stato, età e livello, il profilo deterministico di chi porta il sensore con la distanza d'avviso, il segnale di affidabilità persa e il dominio di validità (\S\ref{sec:lim}).
\end{enumerate}

\subsection{Sintesi dei risultati e livello di prova}\label{sec:intro-sintesi}

La Tabella~\ref{tab:intro-sintesi} riporta i risultati principali. Ogni risultato porta il livello di prova: \prova{D} dimostrato nel modello dichiarato; \prova{C} verificato per calcolo; \prova{S} verificato in simulazione; \prova{P} misurato su dati pubblici; \prova{N} non verificato sperimentalmente. Le simulazioni sono di due tipi: \emph{ideale}, cioè con modello ideale del sensore (posizioni al suolo con rumore gaussiano di $5\unita{cm}$, probabilità di rilevamento per fotogramma $P_D$, sensore e ostacolo fermi, un soggetto per ombra), e \emph{da scheda}, cioè su profondità simulata dalle schede tecniche di tre camere a stereo attivo a $424\times240$ pixel, a posa nota. Gli intervalli fra parentesi quadre sono di Wilson al 95\,\%.

\begin{table}[p]
\centering\footnotesize
\setlength{\tabcolsep}{4pt}
\caption{Risultati principali, con condizioni, numerosità e livello di prova.}\label{tab:intro-sintesi}
\begin{tabularx}{\textwidth}{@{}>{\RaggedRight\arraybackslash}p{2.8cm}>{\RaggedRight\arraybackslash}p{4.1cm}L>{\RaggedRight\arraybackslash}p{1.2cm}>{\RaggedRight\arraybackslash}p{1.05cm}@{}}
\toprule
Risultato & Condizioni & Valore & Prova & Dove\\
\midrule
Nessuna cella non osservata è dichiarata «libero» & regola di scrittura unica; per l'inverso: profondità valida corretta, posa esatta, suolo piano & proposizione; fuori dalle ipotesi (vetro, errore di posa di $1^\circ$) cade & \prova{D} & \S\ref{sec:map}\\
«Libero» della cella & ideale, errore d'inclinazione entro $0{,}5^\circ$; 300 condizioni per 30 fotogrammi, due lanci & 0 celle «libero» su un ostacolo; per condizione $[0;\,1{,}3\,\%]$ & \prova{S} & \S\ref{sec:map}\\
«Libero» nell'ombra vera & da scheda, livello di rumore nominale; 131 sequenze & 0 su $2\,907\,032$ celle\,$\cdot$\,fotogramma; per sequenza $[0;\,2{,}9\,\%]$ & \prova{S} & \S\ref{sec:val}\\
Contenimento in $\Pres(t)$ e riduzione solo per osservazione & profondità valida corretta, posa entro tolleranza, velocità al più $\vmax$ & proposizioni & \prova{D} & \S\ref{sec:occ}\\
Fascia di guardia di $0{,}40\unita{m}$ & età del dato $0{,}133\unita{s}$, latenza stimata $0{,}155\unita{s}$, $\Delta=0{,}067\unita{s}$; ostacolo fermo & copre $\vmax=1{,}13\unita{m/s}$; $0{,}56\unita{m/s}$ se il bordo dell'ombra sbaglia di $0{,}20\unita{m}$ & \prova{D}\prova{C} & \S\ref{sec:occ}\\
Continuità & ideale; 12 combinazioni per 5 comportamenti; premessa: ingresso nell'ombra riconosciuto & 0 passaggi con un ciclo senza regione su $23\,402$ (6000 camminatori distinti: al più $6{,}4$ su $10\,000$); premessa mancata in 24 su $23\,402$. Con corpo puntiforme: 0 su $24\,835$ $[0;\,0{,}016\,\%]$ & \prova{D}\prova{S} & \S\ref{sec:occ}, \S\ref{sec:val}\\
Copertura di $\Reg_\alpha$ al livello $0{,}90$ & ideale; ostacolo largo $2\unita{m}$; corpo esteso; soggetto sul proprio percorso; $P_D=0{,}9$; bordo dell'ombra noto & $\ell^->0{,}90$ fino a $2\unita{s}$ e $>0{,}85$ a $3\unita{s}$ in 30 combinazioni su 30; peggiori $2125/2279$ ($\ell^-=0{,}921$) e $1754/1922$ ($\ell^-=0{,}899$); da 608 a 2359 casi per combinazione & \prova{S} & \S\ref{sec:val}\\
\quad pannello largo $1\unita{m}$ & come sopra & a $1\unita{s}$ $942/1045$ ($\ell^-=0{,}882$; $0{,}87$ su cinque combinazioni insieme); oltre $1\unita{s}$ nessuna cifra. Con corpo puntiforme: $826/879$ a $2\unita{s}$ ($\ell^-=0{,}922$) & \prova{S} & \S\ref{sec:val}\\
\quad fuori dal dominio & ideale; chi svolta o esce dal lato lontano; bordo dell'ombra stimato più stretto di 10--20\,cm & a 2 e $3\unita{s}$: $0{,}13$--$0{,}70$ con corpo puntiforme ($N$ da 538 a 1791), $0{,}57$--$0{,}79$ con corpo esteso e ostacolo largo $2\unita{m}$; bordo stimato: $0{,}30$--$0{,}46$ a $1\unita{s}$. Resta la sola area d'avviso & \prova{S} & \S\ref{sec:val}\\
Avviso di possibile presenza & ideale; un soggetto per ombra & soggetto nascosto dentro l'area d'avviso per costruzione; avviso non acceso in 29 passaggi su $95\,399$ (ombra esatta) e in 15 su 2000 (ombra stimata più stretta di $0{,}20\unita{m}$) & \prova{D}\prova{S} & \S\ref{sec:occ}\\
Identità alla ricomparsa & ideale, sul percorso, tre disposizioni & $10\,711/10\,757$; $9228/9240$; $6772/6967$ ($\ell^-=0{,}994$; $0{,}998$; $0{,}968$) & \prova{S} & \S\ref{sec:sti}\\
\quad su profondità simulata & da scheda; parametri di partenza; poi gruppi e riaggancio modificati, soggetto a $0{,}5\unita{m}$ dietro il pannello & 0 su 90 $[0;\,4{,}1\,\%]$; poi 176 su 180 $[94;\,99\,\%]$. Un secondo soggetto nella stessa ombra riceve l'identità del primo (19 casi su 20); l'avviso resta & \prova{S} & \S\ref{sec:val}\\
Modello di moto & traiettorie pubbliche di pedoni \cite{thor}; 219 tratti, 9 persone; $q=0{,}2\unita{m^2/s^3}$ & copertura della regione al 90\,\% del filtro: $0{,}977$; $1{,}000$; $1{,}000$ a 1, 2, $3\unita{s}$ ($[0{,}948;\,0{,}990]$ a $1\unita{s}$, tratti contati come indipendenti) & \prova{P} & \S\ref{sec:val}\\
Regola «mobile» & ideale, $0{,}5\unita{m/s}$ da $2\unita{m}$, $N=2000$; da scheda, pannello fermo, 90 passaggi & «in movimento» entro $1\unita{s}$: 100\,\% ($\ell^-=99{,}81\,\%$); pannello fermo dichiarato mobile in 79 su 90 $[79{,}4;\,93{,}0\,\%]$, in 5 su 90 con la soglia adattata alla risoluzione & \prova{S} & \S\ref{sec:sti}\\
Ostacoli fermi & da scheda; una sequenza per caso & entro $15\unita{cm}$ dalla posizione vera: cassa alta $1\unita{m}$ a $1{,}68\unita{m}$, 84 fotogrammi su 84; scatola di $15\unita{cm}$ a $2{,}0\unita{m}$, 10 su 10; la stessa a $0{,}8\unita{m}$ e una lastra di vetro, 0 su 10. Descrivono un caso, non una frequenza & \prova{S} & \S\ref{sec:val}\\
Osservato distinto dal previsto & da scheda; parametri di partenza, poi gruppi e riaggancio modificati & uscita «osservato» con il soggetto nascosto in 8 passaggi su 90, poi in 1 su 140: la soglia 0 non è rispettata & \prova{S} & \S\ref{sec:val}\\
Massa e accelerazione & punto materiale con forza ignota; misure di posizione, $\sigma_p=3\unita{cm}$ assunto, $15\unita{Hz}$ & massa non identificabile; scarto tipo dell'accelerazione $16{,}5\unita{m/s^2}$ contro $0{,}67\unita{m/s^2}$ di valore efficace & \prova{D}\prova{C} & \S\ref{sec:sti}\\
Costo per ciclo & da scheda, 9152 cicli, un filo di CPU da portatile; una registrazione pubblica affollata, 320 cicli & media $6{,}35\unita{ms}$, 95º percentile $18{,}6\unita{ms}$, 4 cicli oltre $\Delta$; media $81\unita{ms}$ sulla registrazione pubblica & \prova{S}\prova{P} & \S\ref{sec:cos}\\
Latenza & somma delle medie delle voci note, nessuna misurata sulla catena & $155\unita{ms}$; il 95º percentile non è disponibile; il tempo dal primo fotogramma all'uscita sonora non è misurato & \prova{C}\prova{N} & \S\ref{sec:lim}\\
Affidabilità persa & simulazione, rumore dell'unità inerziale ipotizzato & segnale in $0{,}13\unita{s}$; 0 falsi in $60\unita{s}$ & \prova{S} & \S\ref{sec:lim}\\
Classe e avviso & per costruzione; sequenze casuali di etichette sulla sola fascia di guardia & 0 riduzioni dell'avviso su $20\,000$, e 0 su altre $20\,000$; mai eseguito sulla catena completa & \prova{D}\prova{S}\prova{N} & \S\ref{sec:occ}\\
Avviso tempestivo per classe & portata $3\unita{m}$, avvicinamento frontale, reazione di $1{,}0\unita{s}$ ipotizzata & fino a $1{,}44\unita{m/s}$: 3 classi su 15 alla velocità tipica, 0 su 15 alla massima & \prova{C} & \S\ref{sec:occ}, \S\ref{sec:lim}\\
\bottomrule
\end{tabularx}
\end{table}

\subsection{Che cosa l'allegato non dimostra}\label{sec:intro-limiti}

\begin{riquadro}{Limiti di validità}
\textbf{Nessuna misura su un sensore fisico.} Ogni risultato è una dimostrazione in un modello dichiarato, un calcolo, una simulazione (con modello ideale del sensore o su profondità simulata dalle schede tecniche) o un'elaborazione di dati pubblici. Copertura della regione, errore del bordo dell'ombra, identità alla ricomparsa e latenza non sono verificati su un sensore \prova{N}. Il simulatore da scheda produce meno pixel senza dato delle registrazioni pubbliche con cui è stato confrontato ($7{,}8\,\%$ al livello di rumore più alto, contro 24--34\,\%): le cifre al livello nominale sono ottimistiche.

\textbf{Dominio.} Camera ferma, a posa nota entro $0{,}5^\circ$; al chiuso; fino a $3\unita{m}$; suolo piano; ostacolo fermo, opaco, che restituisce profondità e nasconde il soggetto per intero; un soggetto per ombra, visto per almeno 4 fotogrammi prima di scomparire. Il livello della regione vale con il bordo dell'ombra noto e per chi resta sul proprio percorso; non è dichiarata alcuna cifra oltre $3\unita{s}$.

\textbf{Classificazione.} Il classificatore è un componente richiesto dal sistema; in questo lavoro non è realizzato né valutato: è simulato con etichette assegnate a mano, eventualmente con errori iniettati per misurare la sensibilità. Nessuna classe diversa da «persona» è stata simulata nella catena completa. La sicurezza delle indicazioni non dipende dall'etichetta.

Il sistema non è un dispositivo di protezione e non sostituisce il bastone, il cane guida o le procedure di sicurezza di un impianto.
\end{riquadro}

\subsection{Ambiti d'uso}\label{sec:intro-ambiti}

La traccia ammette due applicazioni. La prima è il monitoraggio di attività e oggetti in ambienti aerospaziali e industriali, con il sensore su una stazione fissa o su un veicolo senza equipaggio fermo, e un'uscita a schermo e su registro degli eventi. La seconda è l'avviso di presenza per una persona non vedente ferma, con il sensore su un supporto accanto alla persona e un'uscita sonora. Hanno in comune la catena, la mappa, il trattamento dell'occlusione e il principio di sicurezza; cambiano il profilo di chi porta il sensore, il canale d'uscita e il comportamento alla perdita di affidabilità (\S\ref{sec:lim}). In entrambe il sistema svolge solo osservazione, monitoraggio e avviso: nessuna prova è stata fatta con il sensore in moto o indossato.

\subsection{Struttura del documento}\label{sec:intro-struttura}

\S\ref{sec:geo} descrive il sensore e la geometria d'osservazione; \S\ref{sec:map} la rappresentazione dello spazio e la sicurezza del «libero»; \S\ref{sec:sti} il rilevamento, la classificazione, l'associazione e la stima dello stato; \S\ref{sec:occ} l'occlusione; \S\ref{sec:val} il metodo di valutazione e i risultati; \S\ref{sec:cos} il costo di calcolo, le alternative considerate e i lavori correlati; \S\ref{sec:lim} l'uscita, il profilo di chi porta il sensore, la perdita di affidabilità e i limiti. L'Appendice~\ref{sec:par} raccoglie i parametri e l'Appendice~\ref{sec:rip} le indicazioni per rifare i conti.
